El aumento de aplicaciones web que requieren interacción en tiempo real ha impulsado la consolidación de arquitecturas y patrones de comunicación que priorizan la fluidez y la capacidad de respuesta del sistema frente a los enfoques de carga de páginas completas. En plataformas donde conviven simultáneamente un editor de código, una terminal de ejecución y un canal de diálogo con un tutor basado en IA, esta exigencia se vuelve aún más pronunciada, ya que la calidad pedagógica de la experiencia depende directamente de la latencia y la continuidad del flujo de datos entre el cliente y el servidor.